% ==================================================================
%  preamble.tex : 『素粒子論のための代数トポロジー演習』共通プリアンブル
%  ビルド：upLaTeX (jsbook) + dvipdfmx。本体は test.tex。
% ==================================================================
\documentclass[uplatex, dvipdfmx, 11pt]{jsbook}

\usepackage[utf8]{inputenc}
\usepackage[margin=30truemm]{geometry}
% --- 数学 ---
\usepackage{amsmath,amssymb,amsthm,amsfonts}
\usepackage{mathtools}
\usepackage{physics}      % \tr, \Tr, \dd, \pdv など
\usepackage{bm}
\usepackage{tikz-cd}      % 可換図式
\usepackage{slashed}      % ディラック作用素 \slashed{D}
\numberwithin{equation}{chapter}
% --- 図・色 ---
\usepackage{graphicx, color}
\usepackage{tikz}
\usetikzlibrary{intersections,calc,arrows.meta,decorations.markings,decorations.pathreplacing,patterns,3d}
\usepackage{float}
\usepackage{ascmac}
\usepackage{environ}
\usepackage{showkeys}
\parindent = 0pt

% ------------------------------------------------------------------
% 図用のスタイル（辺の同一視の矢印など）
% ------------------------------------------------------------------
\tikzset{
  mid arrow/.style={postaction={decorate},decoration={markings,
      mark=at position .55 with {\arrow{Stealth[length=2.2mm]}}}},
  mid arrows/.style={postaction={decorate},decoration={markings,
      mark=at position .48 with {\arrow{Stealth[length=2.2mm]}},
      mark=at position .60 with {\arrow{Stealth[length=2.2mm]}}}},
  vtx/.style={circle,fill,inner sep=1.1pt},
  hvtx/.style={circle,draw,fill=white,inner sep=1.3pt},
}
% 図の下に付ける説明（solution 環境の中でも使えるよう float は使わない）
\newcommand{\zucap}[1]{\par\nopagebreak{\centering\footnotesize #1\par}\medskip}

% ------------------------------------------------------------------
% 用語（将来の英語化に備えて英語を併記する）
% ------------------------------------------------------------------
\newcommand{\term}[2]{\textbf{#1}（#2）}   % 初出の用語：太字＋（English）
\newcommand{\en}[1]{（#1）}                % 補足の英語

% 計算過程などを書く注釈
\newenvironment{note}{\vspace{1ex}\par\footnotesize\noindent\textbf{【注】}}{\par\vspace{1ex}}

% ------------------------------------------------------------------
% 記号
% ------------------------------------------------------------------
\newcommand{\id}{\mathrm{id}}
\newcommand{\Z}{\mathbb{Z}}
\newcommand{\R}{\mathbb{R}}
\newcommand{\C}{\mathbb{C}}
\newcommand{\Q}{\mathbb{Q}}
\newcommand{\HH}{\mathbb{H}}               % 四元数
\newcommand{\RP}{\mathbb{RP}}
\newcommand{\CP}{\mathbb{CP}}
\DeclareMathOperator{\Hom}{Hom}
\DeclareMathOperator{\Ext}{Ext}
\DeclareMathOperator{\Tor}{Tor}
\DeclareMathOperator{\coker}{coker}
\DeclareMathOperator{\ind}{ind}
\DeclareMathOperator{\sign}{sign}
\DeclareMathOperator{\ch}{ch}
\DeclareMathOperator{\Td}{Td}
\DeclareMathOperator{\Pf}{Pf}
\DeclareMathOperator{\Aut}{Aut}
\DeclareMathOperator{\End}{End}
\DeclareMathOperator{\Ad}{Ad}
\DeclareMathOperator{\ad}{ad}
\DeclareMathOperator{\Spin}{Spin}
\DeclareMathOperator{\Pin}{Pin}
\newcommand{\Ahat}{\hat{A}}

% ------------------------------------------------------------------
% ハイパーリンク
% ------------------------------------------------------------------
\usepackage[dvipdfmx,bookmarks=true,bookmarksnumbered=true,bookmarkstype=toc,
  colorlinks=true,linkcolor=blue,citecolor=blue,urlcolor=blue]{hyperref}
\usepackage{pxjahyper}

% ------------------------------------------------------------------
% 定理系の囲み（tcolorbox）
%   使い方： \begin{mythm}{タイトル}{ラベル} ... \end{mythm}
%   ラベルは thm:ラベル, def:ラベル, ... として参照できる
% ------------------------------------------------------------------
\usepackage[many]{tcolorbox}
\tcbset{
  myboxstyle/.style={
    enhanced, breakable,
    colback=white, coltitle=white, fonttitle=\bfseries,
    boxrule=0.8pt, arc=1.5mm,
    attach boxed title to top left={xshift=3mm, yshift=-2mm},
    top=3mm, before skip=3mm, after skip=3mm,
    description font=\normalfont\bfseries, separator sign={：}
  }
}
\newtcbtheorem[number within=chapter]{mythm}{定理}{myboxstyle,
  colframe=red!75!black,
  boxed title style={colback=red!50,colframe=red!75!black,boxrule=0.8pt,arc=1.5mm}}{thm}
\newtcbtheorem[use counter from=mythm]{mydef}{定義}{myboxstyle,
  colframe=blue!70!black,
  boxed title style={colback=blue!60!gray!80,colframe=blue!70!black,boxrule=0.8pt,arc=1.5mm}}{def}
\newtcbtheorem[use counter from=mythm]{mylem}{補題}{myboxstyle,
  colframe=green!50!black,
  boxed title style={colback=green!50!gray,colframe=green!50!black,boxrule=0.8pt,arc=1.5mm}}{lem}
\newtcbtheorem[use counter from=mythm]{myprop}{命題}{myboxstyle,
  colframe=orange!80!black,
  boxed title style={colback=orange!80!black,colframe=orange!80!black,boxrule=0.8pt,arc=1.5mm}}{prop}
\newtcbtheorem[use counter from=mythm]{myex}{例}{myboxstyle,
  colframe=gray!70!black,
  boxed title style={colback=gray!70!black,colframe=gray!70!black,boxrule=0.8pt,arc=1.5mm}}{ex}
% 物理との対応をまとめる囲み（番号なし）
\newtcolorbox{physbox}[1][]{enhanced, breakable, colback=yellow!4, colframe=orange!60!black,
  boxrule=0.6pt, arc=1.5mm, fonttitle=\bfseries, title={物理との対応#1}}
% 旧版互換：章冒頭の「この章で用いる事項」
\newtcolorbox{summary}{breakable, enhanced, colback=gray!5, colframe=gray!60!black,
  boxrule=0.6pt, arc=1.5mm, fonttitle=\bfseries, title=この章で用いる事項}

% ------------------------------------------------------------------
% 問題と解答
%   problem 環境は本文に置き、solution 環境の中身は \printsolutions で巻末にまとめて出力する
%   注意：solution の中では # と verbatim を使わない（中身をマクロに蓄えるため）
% ------------------------------------------------------------------
\theoremstyle{definition}
\newtheorem{problem}{問題}[chapter]
\renewcommand{\proofname}{証明}
\newcommand{\allsolutions}{}
\newcommand{\currentchaptertitle}{}
\let\origchaptermark\chaptermark
\renewcommand{\chaptermark}[1]{\gdef\currentchaptertitle{#1}\origchaptermark{#1}}
\newcounter{solchapter}
\NewEnviron{solution}{%
  \ifnum\value{chapter}>\value{solchapter}%
    \setcounter{solchapter}{\value{chapter}}%
    \xdef\allsolutions{\unexpanded\expandafter{\allsolutions}%
      \noexpand\solutionchapter{\thechapter}{\unexpanded\expandafter{\currentchaptertitle}}}%
  \fi
  \xdef\allsolutions{\unexpanded\expandafter{\allsolutions}%
    \noexpand\solutionitem{\theproblem}{\unexpanded\expandafter{\BODY}}}%
}
\newcommand{\solutionchapter}[2]{\section*{第#1章\quad #2}}
\newcommand{\solutionitem}[2]{\par\medskip\noindent\textbf{問題 #1 の解答}\par\nopagebreak #2\qed\par}
\newcommand{\printsolutions}{%
  \chapter*{解答}
  \addcontentsline{toc}{chapter}{解答}
  \markboth{解答}{解答}
  \allsolutions
}
